% Sources for Auto-generated glossary entries. Order preserved via sort= (original hand-curated order).

\newglossaryentry{rbf}{name={\textsf{\textcolor{red}{\texttt{.rbf}}}},sort={0001},description={Raw Binary File. \term{FPGA} configuration image produced by \term{Quartus Prime} from a \texttt{.sof} file; this is the format the \term{HPS} \term{FPGA Manager} loads at run time to configure the fabric. \cite{CV-TRM}}}
\newglossaryentry{sof}{name={\textsf{\texttt{.sof}}},sort={0002},description={SRAM Object File. \term{Quartus Prime}'s native \term{FPGA} configuration file, used to program the device over \term{JTAG} and as the source from which an \texttt{.rbf} is converted.}}
\newglossaryentry{altera}{name={\textsf{\textcolor{red}{\term{Altera}}}},sort={0003},description={\term{FPGA} vendor that developed the \term{Cyclone} device family and the \term{Quartus} design suite. Acquired by \term{Intel} in 2015 and re-established as an independent \term{Intel} subsidiary in 2024; both names appear in the tools and reference manuals cited here. \citegloss{ALTERA-COMPANY}}}
\newglossaryentry{apt}{name={\textsf{\textcolor{red}{\texttt{apt}}}},sort={0004},description={Package manager of Debian-based \term{Linux} distributions, used on the development host to install build dependencies. \citegloss{APT-UBUNTU-MANPAGE}}}
\newglossaryentry{armbaremetaltoolchain}{name={\textsf{\textcolor{red}{ARM bare-metal toolchain}}},sort={0005},description={Cross-compiler suite targeting an \term{ARM} processor with no operating system (\texttt{\seqsplit{arm-none-eabi-*}}). It provides no system calls and links against a minimal \term{C} runtime; used to build the \prjname{FPGAsteel} projects. \cite{XPACK-ARM-GCC}}}
\newglossaryentry{armgnutoolchain}{name={\textsf{\textcolor{red}{ARM GNU toolchain}}},sort={0006},description={Cross-compiler suite targeting \term{ARM} \term{Linux} (\texttt{\seqsplit{arm-linux-gnueabihf-*}}). It builds kernel and user-space code for the \term{Cortex-A9} of the \term{HPS} on the development host. \cite{ARM-GNU-TOOLCHAIN}}}
\newglossaryentry{avalonmm}{name={\textsf{\textcolor{red}{\term{Avalon-MM} (Avalon Memory-Mapped)}}},sort={0007},description={\term{Intel} address-based interconnect interface: a master issues read and write transactions into the address space of a slave. It carries control and status accesses to the custom cores. \cite{AVALON-SPEC}}}
\newglossaryentry{avalonst}{name={\textsf{\textcolor{red}{\term{Avalon-ST} (Avalon Streaming)}}},sort={0008},description={\term{Intel} unidirectional point-to-point interface carrying a continuous stream of data from a source to a sink, with flow control and optional packet delimiters. Data moves one beat at a time, each beat carrying an integer number of symbols. The sink is the receiving endpoint. \cite{AVALON-SPEC}}}
\newglossaryentry{backpressure}{name={\textsf{\textcolor{red}{backpressure}}},sort={0009},description={Flow-control mechanism of \term{Avalon-ST}: the sink deasserts its ready signal to stall the source when it cannot accept further data, and the stall propagates upstream through the pipeline. \cite{AVALON-SPEC}}}
\newglossaryentry{balenaetcher}{name={\textsf{\textcolor{red}{balenaEtcher}}},sort={0010},description={Cross-platform graphical utility for writing a disk image to an \term{SD} card or \term{USB} drive. \citegloss{BALENAETCHER-SITE}}}
\newglossaryentry{baremetal}{name={\textsf{bare metal}},sort={0011},description={Execution model in which the application runs directly on the processor with no operating system underneath, taking full control of memory, interrupts and peripherals. It is the model adopted by \prjname{FPGAsteel}.}}
\newglossaryentry{bayermosaic}{name={\textsf{\textcolor{red}{Bayer mosaic}}},sort={0012},description={Color filter array in which every pixel of the image records only one of the three primaries (i.e.\ Red, Green, Blue), arranged in a repeating 2$\times$2 pattern with two green elements. A full-color image can be reconstructed from it by demosaicing. \citegloss{US-PATENT-3971065}}}
\newglossaryentry{beat}{name={\textsf{\textcolor{red}{beat}}},sort={0013},description={Unit of transfer of an \term{Avalon-ST} link: a single clock cycle in which the source presents valid data and the sink accepts it. The amount of data moved per beat is fixed by the interface and is equal to the number of symbols per beat multiplied by the symbol width: an interface carrying four 8-bit symbols per beat moves 32 bits, that is four bytes, on every beat. \cite{AVALON-SPEC}}}
\newglossaryentry{bootchain}{name={\textsf{boot chain}},sort={0014},description={Ordered sequence of programs that bring the system from power-on to the running application/operating system.}}
\newglossaryentry{bootrom}{name={\textsf{\textcolor{red}{Boot \term{ROM}}}},sort={0015},description={Fixed-function code embedded in the \term{HPS} silicon at fabrication time; the first code executed on reset, and immutable thereafter. It samples the boot source selection pins, loads \term{SPL} from the selected source into on-chip \term{RAM}, and transfers control to \term{SPL}. \cite{CV-TRM}}}
\newglossaryentry{bootscript}{name={\textsf{\textcolor{red}{boot script}}},sort={0016},description={Text file of \term{U-Boot} commands. Once compiled into \texttt{u-boot.scr}, it is executed automatically by \term{U-Boot} at startup. It can be used, for example, to program the \term{FPGA}, load the kernel or the bare-metal application image into memory and hand over control to them. \cite{SUB-UBOOT}}}
\newglossaryentry{bootloader}{name={\textsf{\textcolor{red}{bootloader}}},sort={0017},description={Program that initializes the processor, the external memory and the essential peripherals, then loads the next stage of the boot chain. Here it is \term{U-Boot} together with its \term{SPL}. \cite{SUB-UBOOT}}}
\newglossaryentry{bootz}{name={\textsf{\textcolor{red}{\texttt{bootz}}}},sort={0018},description={\term{U-Boot} command that boots a \term{Linux} kernel in \term{zImage} format, given the load addresses of the kernel and of the device tree blob. \cite{SUB-UBOOT}}}
\newglossaryentry{bufferpool}{name={\textsf{buffer pool}},sort={0019},description={Fixed set of identical buffers, allocated once at start-up and then recycled: the producer takes a free one, fills it and passes it to the consumer, which releases it back to the pool when done. The technique generalizes to any resource whose creation is expensive; allocating up front keeps that cost off the critical path, makes timing predictable and puts a hard ceiling on the memory in use. Here the pooled elements are video frame buffers.}}
\newglossaryentry{buildroot}{name={\textsf{\textcolor{red}{\term{Buildroot}}}},sort={0020},description={Build system that generates, from a single configuration, a complete embedded \term{Linux} root filesystem, together with the cross-toolchain that will be used to compile the application that will run on it. \cite{SUB-BUILDROOT}}}
\newglossaryentry{busybox}{name={\textsf{\textcolor{red}{\term{BusyBox}}}},sort={0021},description={Single executable providing compact implementations of the common \term{UNIX} utilities and of the init process; it forms the core of the root filesystem built with \term{Buildroot}. \citegloss{BUSYBOX-SITE}}}
\newglossaryentry{cpp}{name={\textsf{\textcolor{red}{\term{C++}}}},sort={0022},description={Compiled, statically typed programming language with object-oriented and generic facilities; used here for the \texttt{vision\_core} application and for the bare-metal projects. \citegloss{CPPREFERENCE}}}
\newglossaryentry{cache}{name={\textsf{cache}},sort={0023},description={Small, fast associative memory that holds recently accessed lines of main memory, so that repeated accesses are served without paying the latency of external \term{RAM}. The \term{Cortex-A9} cluster has two private \term{L1} caches per core, one for data and one for instructions, and a single unified (i.e.\ instruction and data) \term{L2} cache, shared between the two cores.}}
\newglossaryentry{claudecode}{name={\textsf{\textcolor{red}{\term{Claude Code}}}},sort={0024},description={\term{Anthropic}'s \term{AI} coding assistant, used throughout this project both to author code and to analyse third-party sources, such as the \term{Linux} kernel's \term{mSGDMA} driver (Section~\ref{sec:fpgalix-camera-short-frames}) and its \term{SMP} bring-up sequence (Section~\ref{sec:fpgasteel-dualcore-about}). \citegloss{CLAUDECODE-SITE}}}
\newglossaryentry{cpucore}{name={\textsf{\term{CPU} core}},sort={0025},description={Independent instruction-executing unit of a processor. The \term{Cyclone V} \term{HPS} integrates a dual-core \term{ARM} \term{Cortex-A9}.}}
\newglossaryentry{csr}{name={\textsf{\term{CSR} (Control and Status Register)}},sort={0026},description={Register block through which software configures a hardware core and reads back its state. For example, the \term{mSGDMA} dispatcher exposes one such block.}}
\newglossaryentry{cyclonevsoc}{name={\textsf{\textcolor{red}{\term{Cyclone V SoC}}}},sort={0027},description={\term{Intel}/\term{Altera} device family that combines an \term{FPGA} fabric and a hard processor system on the same die. The \term{Terasic} \term{DE1-SoC} board used in this work carries one. \cite{CV-TRM}}}
\newglossaryentry{daemon}{name={\textsf{daemon}},sort={0028},description={Background process of a \term{UNIX}-like system, detached from any terminal and normally started at boot. \texttt{hexip\_daemon} is deployed in this form.}}
\newglossaryentry{dd}{name={\textsf{\textcolor{red}{\texttt{dd}}}},sort={0029},description={\term{UNIX} command that copies data block by block between files and devices; used here to write the \term{SD} card image and to place the bootloader in its raw partition. \citegloss{DD-MANPAGE}}}
\newglossaryentry{demosaicing}{name={\textsf{\textcolor{red}{demosaicing}}},sort={0030},description={Interpolation that reconstructs a full set of color components for every pixel from the single-primary samples of a \term{Bayer mosaic}. \citegloss{US-PATENT-5629734}}}
\newglossaryentry{devicetree}{name={\textsf{\textcolor{red}{Device Tree (\texttt{DTS}, \texttt{.dtsi}, \texttt{DTB})}}},sort={0031},description={Description of the hardware that cannot be discovered at run time, passed to the \term{Linux} kernel at boot. On an \term{SoC} most peripherals sit at fixed addresses and cannot be enumerated over an auto-discovering bus, so the kernel must be told what exists, where its registers are, which interrupt and clock it uses and which driver should bind to it. That description is a tree of nodes, one per device, written as a Device Tree Source file (\texttt{.dts}) that may include shared fragments (\texttt{.dtsi}), typically one describing the \term{SoC} family and one the specific board, and compiled by the device tree compiler into a binary Device Tree Blob (\texttt{.dtb}), which the bootloader loads into memory and hands to the kernel. Keeping it separate from the kernel image means the same kernel can serve different boards. \citegloss{DEVICETREE-ORG}}}
\newglossaryentry{dhcpcd}{name={\textsf{\textcolor{red}{\texttt{dhcpcd}}}},sort={0032},description={\term{DHCP} client daemon. \term{DHCP} (Dynamic Host Configuration Protocol) is the protocol by which a machine joining a network asks a server for the parameters it needs to communicate: its own \term{IP} address, the netmask, the default gateway and the \term{DNS} servers. The address is granted for a limited time, the lease, and is renewed periodically. On the board \texttt{dhcpcd} performs this exchange at boot, so that the system is reachable over the network without any address being configured by hand. \citegloss{RFC2131-DHCP,DHCPCD-SITE}}}
\newglossaryentry{die}{name={\textsf{die}},sort={0033},description={Single piece of silicon on which an integrated circuit is fabricated. In the \term{Cyclone V SoC} the \term{FPGA} fabric and the \term{HPS} share one die.}}
\newglossaryentry{diskimagefile}{name={\textsf{disk image file}},sort={0034},description={Regular file holding the byte-for-byte contents of a storage device, partition table included; writing it to an \term{SD} card reproduces the original medium exactly.}}
\newglossaryentry{dma}{name={\textsf{\textcolor{red}{\term{DMA} (Direct Memory Access)}}},sort={0035},description={Transfer of data between a peripheral and memory carried out by a dedicated engine instead of by the processor, which is therefore free to do other work. \citegloss{BRYANT-OHALLARON-CSAPP}}}
\newglossaryentry{eclipseembeddedcdt}{name={\textsf{\textcolor{red}{\term{Eclipse Embedded CDT}}}},sort={0036},description={Ready-made distribution of the \term{Eclipse} \term{IDE} (integrated development environment) aimed at embedded \term{C}/\term{C++} work, downloaded already complete with the plug-ins it needs, so that neither a generic \term{Eclipse} nor any manual plug-in installation is involved. It can handle the \term{ARM} toolchains and provides a graphical front end to \term{GDB} and \term{OpenOCD}. In this project it is used only as a debugger for the bare-metal applications, driving breakpoints, stepping and memory inspection on the target from the same window as the source code; the builds are run with \texttt{make} from the command line. \cite{ECLIPSE-EMBED-CDT}}}
\newglossaryentry{environmentvariables}{name={\textsf{environment variables}},sort={0037},description={Named values held in the environment of a shell and inherited by the processes it launches. The most important of them is \texttt{PATH}. When a command is typed without a full path, the shell does not search the whole disk: it walks the list of directories held in \texttt{PATH}, in order, and runs the first executable with that name it finds. A tool that is installed but whose directory is not listed there is therefore reported as not found, and if two versions are installed the one in the earlier directory wins.}}
\newglossaryentry{exception}{name={\textsf{\textcolor{red}{exception (\term{C++})}}},sort={0038},description={Language mechanism that transfers control from the point where an error is signaled to a matching handler, unwinding the stack in between. \citegloss{GREGOIRE-PROFESSIONAL-CPP}}}
\newglossaryentry{ffmpeg}{name={\textsf{\textcolor{red}{\texttt{ffmpeg}}}},sort={0039},description={Command-line tool of the \term{FFmpeg} project for decoding, encoding, transcoding and streaming audio and video. \citegloss{FFMPEG-SITE}}}
\newglossaryentry{ffplay}{name={\textsf{\textcolor{red}{\texttt{ffplay}}}},sort={0040},description={Minimal media player built on the \term{FFmpeg} libraries, used here to display the video stream produced by the board. \citegloss{FFMPEG-SITE}}}
\newglossaryentry{fpga}{name={\textsf{\term{FPGA} (Field-Programmable Gate Array)}},sort={0041},description={Integrated circuit whose logic resources and interconnect are configured after manufacturing, so that arbitrary digital hardware can be implemented and later modified.}}
\newglossaryentry{fpgabitstream}{name={\textsf{\textcolor{red}{\term{FPGA} bitstream}}},sort={0042},description={Configuration data defining the logic implemented in the \term{FPGA} fabric. \term{Quartus Prime} generates it as a \texttt{.sof} file, which is converted to \texttt{.rbf} for loading by the \term{HPS}. \cite{ALTERA-QUARTUS}}}
\newglossaryentry{fpgamanager}{name={\textsf{\textcolor{red}{\term{FPGA Manager}}}},sort={0043},description={Hardware block of the \term{Cyclone V} through which the \term{HPS} configures the \term{FPGA} fabric at run time from an \texttt{.rbf} image. Any software running on the \term{HPS} can drive it: here it is \term{U-Boot}, from the boot script, that configures the fabric before handing control to the kernel or to the bare-metal application. \term{Linux} also exposes it as a subsystem, so that the fabric can be reconfigured later from the running system. \cite{CV-TRM}}}
\newglossaryentry{fpgalix}{name={\textsf{\prjname{FPGAlix}}},sort={0044},description={Embedded-\term{Linux} branch of this work: \term{OV7670} acquisition in the \term{FPGA} feeding a \term{V4L2} kernel driver and user-space video applications running under \term{Linux}.}}
\newglossaryentry{fpgasteel}{name={\textsf{\prjname{FPGAsteel}}},sort={0045},description={Bare-metal branch of this work: the same acquisition hardware as \prjname{FPGAlix}, extended with a \term{VGA} output path and driven directly by the \term{Cortex-A9} with no operating system.}}
\newglossaryentry{fsm}{name={\textsf{\term{FSM} (Finite State Machine)}},sort={0046},description={Sequential logic construct with a finite set of states and defined transitions between them; the common way of describing control logic.}}
\newglossaryentry{gdb}{name={\textsf{\textcolor{red}{\texttt{gdb}}}},sort={0047},description={\term{GNU} Debugger. Controls the execution of a program, inspects memory and registers, and here connects remotely to \texttt{gdbserver} on the target or to \term{OpenOCD} through \term{JTAG}. \citegloss{GDB-SITE}}}
\newglossaryentry{gdbserver}{name={\textsf{\textcolor{red}{\texttt{gdbserver}}}},sort={0048},description={Small agent running on the target that executes the program under the control of a remote \texttt{gdb} over a network or serial link. \citegloss{GDB-SITE}}}
\newglossaryentry{ghrd}{name={\textsf{\textcolor{red}{\term{GHRD} (Golden Hardware Reference Design)}}},sort={0049},description={Reference \term{Quartus} project supplied by \term{Intel}/\term{Altera} for a given \term{SoC} board, used as the starting point for a custom design. \cite{GHRD}}}
\newglossaryentry{gic}{name={\textsf{\textcolor{red}{\term{GIC} (Generic Interrupt Controller)}}},sort={0050},description={\term{ARM} interrupt controller that collects the interrupt sources, prioritizes them and distributes them to the \term{CPU} cores. \cite{CV-TRM}}}
\newglossaryentry{git}{name={\textsf{\textcolor{red}{\texttt{git}}}},sort={0051},description={Version control system, widely used in software development: it records the successive states of a set of files, so that any earlier version can be recovered and every change traced back to its author. It is called distributed because each copy of a repository carries the whole history and works on its own, without a central server. The two repositories of this work are published on \term{GitHub}. \citegloss{GIT-SITE}}}
\newglossaryentry{github}{name={\textsf{\textcolor{red}{\term{GitHub}}}},sort={0052},description={Web service that hosts \texttt{git} repositories online, so that they can be shared and cloned, and adds a web interface, issue tracking and access control. Access to the repositories goes through \texttt{git} itself, with the usual commands; the web interface is only an addition. It is the most widespread of several equivalent services. \citegloss{GITHUB-SITE}}}
\newglossaryentry{gpio0header}{name={\textsf{\textcolor{red}{\term{GPIO\_0} header}}},sort={0053},description={Forty-pin expansion connector of the \term{DE1-SoC} board, on which the camera adapter board is plugged. \cite{DE1-SCH}}}
\newglossaryentry{gstreamer}{name={\textsf{\textcolor{red}{\term{GStreamer}}}},sort={0054},description={Multimedia framework that builds an audio or video processing chain by connecting elements in a pipeline, each performing one step: capture, format conversion, encoding, display, transmission over the network. The elements come ready made with the framework, several hundred of them covering the common formats and devices, so that a pipeline is usually a matter of listing the right ones rather than writing code; new custom elements can be added if none fits. \citegloss{GSTREAMER-SITE}}}
\newglossaryentry{gui}{name={\textsf{\term{GUI} (Graphical User Interface)}},sort={0055},description={Interface based on windows, menus and pointer interaction, as opposed to a command-line interface, the mode of interaction a shell provides.}}
\newglossaryentry{h264}{name={\textsf{\textcolor{red}{\term{H.264}}}},sort={0056},description={Video compression standard (\term{ITU-T} \term{H.264}, also known as \term{MPEG-4} \term{AVC}), considerably more efficient than \term{MJPEG} at the cost of a much higher computational load. \citegloss{ITU-T-H264}}}
\newglossaryentry{heap}{name={\textsf{heap and stack}},sort={0057},description={The two dynamic regions of the memory of a program. The heap serves allocations of arbitrary size and lifetime, requested explicitly; the stack holds call frames, local variables and return addresses in last-in first-out order. A large buffer declared inside a function, main included, is a local variable and therefore goes on the stack, whose size is fixed and modest: this is why such buffers are usually declared static or allocated on the heap instead. Global and static data belongs to neither region, being placed in areas of the image (\texttt{.data} and \texttt{.bss}) whose size is fixed at compile time.}}
\newglossaryentry{hps}{name={\textsf{\textcolor{red}{\term{HPS} (Hard Processor System)}}},sort={0058},description={\term{ARM} subsystem of the \term{Cyclone V SoC}: dual-core \term{Cortex-A9}, caches, \term{SDRAM} controller, peripherals and the bridges towards the \term{FPGA} fabric. It is called hard because it is built as fixed silicon, like an ordinary microprocessor, and not out of the configurable resources of the \term{FPGA} fabric: it cannot be modified, but it runs faster, occupies less area and, above all, leaves the whole fabric free for the user design. \cite{CV-TRM}}}
\newglossaryentry{href}{name={\textsf{\textcolor{red}{\term{HREF} (Horizontal REFerence)}}},sort={0059},description={Output of the \term{OV7670} that is asserted while valid pixel data of a line is present on the data bus, thus delimiting the active part of each row. \cite{OV7670-DS}}}
\newglossaryentry{hwlib}{name={\textsf{\textcolor{red}{\texttt{hwlib}}}},sort={0060},description={\term{Intel}/\term{Altera} bare-metal \term{C} library providing register definitions and helper functions for the peripherals of the \term{HPS}. \cite{SUB-HWLIB}}}
\newglossaryentry{i2c}{name={\textsf{\textcolor{red}{\term{I\textsuperscript{2}C} and \term{SCCB}}}},sort={0061},description={Two-wire serial bus, made of a clock line and a bidirectional data line, used to configure peripherals. \citegloss{TI-I2C-GUIDE} \term{SCCB} is the nearly identical \term{OmniVision} variant through which the registers of the \term{OV7670} are written. \cite{OV7670-DS}}}
\newglossaryentry{ipcore}{name={\textsf{\term{IP} core (Intellectual Property core)}},sort={0062},description={A self-contained hardware module with well-defined interfaces that can be instantiated as a single unit within an \term{FPGA} design and reused across multiple projects. \term{IP} cores may originate either from vendor-provided libraries, such as the \term{\seqsplit{mSGDMA}} controller and \term{FIFO} components, or from custom developments specific to a given project, as in the case of the acquisition and video output cores presented in this work. The defining characteristic of an \term{IP} core is therefore its packaging and reusability rather than its origin. The term should not be confused with \term{IP} in the context of the Internet Protocol.}}
\newglossaryentry{irq}{name={\textsf{\term{IRQ} (Interrupt Request)}},sort={0063},description={Signal by which a peripheral asks the processor to suspend the current flow of execution and run a handler.}}
\newglossaryentry{isr}{name={\textsf{\term{ISR} (Interrupt Service Routine)}},sort={0064},description={Handler function the processor invokes in response to an interrupt, suspending whatever was running and returning control to it once the handler completes. In this work, cache and \term{DMA} transfer-complete notifications are delivered this way.}}
\newglossaryentry{jtag}{name={\textsf{\textcolor{red}{\term{JTAG}}}},sort={0065},description={\term{IEEE} 1149.1 serial interface used to program the \term{FPGA}, to access the on-chip debug logic and to control the processor from a host. \citegloss{IEEE-1149-1}}}
\newglossaryentry{kernel}{name={\textsf{\textcolor{red}{kernel (\term{Linux})}}},sort={0066},description={Privileged software layer interposed between the hardware and the programs: it assigns memory and processor time to the processes, arbitrates access to the devices through its drivers, and exposes the system calls through which a program requests these services. \term{Linux} is this layer alone; the shell, the utilities and the libraries that turn it into a usable system belong to the root filesystem and are developed independently of it. \cite{SUB-LINUX}}}
\newglossaryentry{loopdevice}{name={\textsf{loop device}},sort={0067},description={\term{Linux} pseudo-device that makes a regular file accessible as a block device, so that a disk image can be partitioned and mounted as if it were physical media.}}
\newglossaryentry{lwhpsbridge}{name={\textsf{\textcolor{red}{\term{LW HPS bridge} (Lightweight \term{HPS}-to-\term{FPGA} bridge)}}},sort={0068},description={A 32-bit \term{Avalon-MM} bridge that allows the \term{HPS} to access memory-mapped registers in the \term{FPGA} fabric. It is intended for low-bandwidth control and status accesses, not for high-throughput data transfers. \cite{CV-TRM}}}
\newglossaryentry{make}{name={\textsf{\textcolor{red}{\texttt{make}}}},sort={0069},description={Classic \term{UNIX} build automation tool. Through a \texttt{Makefile}, the developer sets the target (what to build), the files it depends on (the dependencies), and the commands needed to build it. When \texttt{make} is invoked, it reads the \texttt{Makefile} and runs those commands to compile and link the program. \citegloss{GNU-MAKE-SITE}}}
\newglossaryentry{menuconfig}{name={\textsf{\texttt{menuconfig}}},sort={0070},description={Text-based configuration front end, invoked as \texttt{make menuconfig}, that sets the build options of a software project. Used here for \term{U-Boot}, the \term{Linux} kernel, and \term{Buildroot}.}}
\newglossaryentry{mjpeg}{name={\textsf{\textcolor{red}{\term{MJPEG} (Motion JPEG)}}},sort={0071},description={Video format in which every frame is coded independently as a \term{JPEG} image. It is simple and low-latency, but far less efficient than \term{H.264}. \citegloss{ITU-T-T81-JPEG}}}
\newglossaryentry{mmap}{name={\textsf{\textcolor{red}{\texttt{mmap}}}},sort={0072},description={\term{POSIX} system call that maps a file or a device directly into the calling process's address space: once mapped, ordinary pointer accesses read and write the underlying file or device directly, without going through \texttt{read()}/\texttt{write()} calls or an intermediate kernel copy. Used here for the \term{V4L2} camera buffers and for the register windows of the \term{FPGA} peripherals. \citegloss{MMAP-MANPAGE}}}
\newglossaryentry{mmu}{name={\textsf{\textcolor{red}{\term{MMU} (Memory Management Unit)}}},sort={0073},description={Hardware unit of the \term{Cortex-A9} that translates the virtual addresses issued by software into physical ones, according to a page table in memory, and enforces the access and caching attributes of each region. \cite{CV-TRM}}}
\newglossaryentry{mpfe}{name={\textsf{\textcolor{red}{\term{MPFE} (Multi-Port Front End)}}},sort={0074},description={Arbitration front end of the \term{Cyclone V} \term{SDRAM} controller, through which the \term{HPS} and the masters implemented in the \term{FPGA} fabric access external \term{SDRAM} (\term{DDR}) memory. \cite{CV-TRM}}}
\newglossaryentry{mpu}{name={\textsf{\textcolor{red}{\term{MPU}}}},sort={0075},description={\term{Cyclone V} \term{HPS} term for the dual-core \term{ARM} \term{Cortex-A9} complex: the two cores together with the \term{L2} cache, as distinguished from the rest of the \term{HPS} (peripherals, bridges, \term{L3} interconnect). \cite{CV-TRM}}}
\newglossaryentry{msgdma}{name={\textsf{\textcolor{red}{\term{mSGDMA} (Modular Scatter-Gather \term{DMA})}}},sort={0076},description={\term{Intel}/\term{Altera} \term{DMA} \term{IP} core that offloads the \term{CPU} from moving data between two memory spaces, or between an \term{Avalon-ST} stream and memory, in one direction or the other depending on configuration. Its modular architecture lets it adapt to each situation: a hardware descriptor prefetcher, for instance, can be attached to the base module so it re-arms itself from on-chip memory without \term{CPU} intervention. \cite{EP-IPUG}}}
\newglossaryentry{neon}{name={\textsf{\textcolor{red}{\term{NEON}}}},sort={0077},description={\term{SIMD} (single instruction, multiple data) extension of the \term{ARM} \term{Cortex-A} architecture, which applies one operation to several data elements at once. Here it is used to accelerate pixel processing. \cite{CV-TRM} \citegloss{ARM-DDI-0409I}}}
\newglossaryentry{ntp}{name={\textsf{\textcolor{red}{\term{NTP} (Network Time Protocol)}}},sort={0078},description={Protocol that synchronizes the system clock over the network, relevant on a board that has no battery-backed real-time clock. \citegloss{RFC5905-NTP}}}
\newglossaryentry{onchipdescriptormemory}{name={\textsf{\textcolor{red}{on-chip descriptor memory}}},sort={0079},description={Small memory implemented in the \term{FPGA} fabric. Here it holds the \term{mSGDMA} prefetcher's descriptors, so that they are fetched without accessing external \term{SDRAM}. \cite{EP-IPUG}}}
\newglossaryentry{onchipram}{name={\textsf{\textcolor{red}{on-chip \term{RAM} (\term{SRAM})}}},sort={0080},description={Used in this work in two distinct senses: a small \term{SRAM} block inside the \term{HPS} itself, used by the boot \term{ROM} and \term{SPL} before the \term{SDRAM} controller is initialized (see \term{SPL}); and memory implemented in the embedded blocks of the \term{FPGA} fabric, offering low and deterministic latency at a far smaller capacity than external \term{SDRAM}. In this second sense, it is the \term{mSGDMA} prefetcher's descriptor memory (see on-chip descriptor memory). \cite{CV-TRM}}}
\newglossaryentry{opencv}{name={\textsf{\textcolor{red}{\term{OpenCV}}}},sort={0081},description={Open-source computer vision library, used in \texttt{vision\_core} for format conversion and image processing. \citegloss{OPENCV-SITE}}}
\newglossaryentry{openocd}{name={\textsf{\textcolor{red}{\term{OpenOCD}}}},sort={0082},description={Open On-Chip Debugger. It bridges a \term{JTAG} probe to \texttt{gdb}, enabling bare-metal debugging of the \term{Cortex-A9}. \cite{XPACK-OPENOCD}}}
\newglossaryentry{openssh}{name={\textsf{\textcolor{red}{\term{OpenSSH}}}},sort={0083},description={Reference implementation of the \term{SSH} protocol, providing the \texttt{sshd} server and the \texttt{ssh} and \texttt{scp} clients. \citegloss{OPENSSH-SITE}}}
\newglossaryentry{ov7670}{name={\textsf{\textcolor{red}{\term{OV7670}}}},sort={0084},description={\term{OmniVision} \term{CMOS} image sensor at \term{VGA} resolution, with an eight-bit parallel data bus, \term{PCLK}, \term{HREF} and \term{VSYNC} timing signals and an \term{I\textsuperscript{2}C}-like (i.e. \term{SCCB}) configuration interface. \cite{OV7670-DS}}}
\newglossaryentry{pcb}{name={\textsf{\term{PCB} (Printed Circuit Board)}},sort={0085},description={Substrate that mechanically supports and electrically interconnects its components; here the transmitter and receiver boards of the camera adapter.}}
\newglossaryentry{pclk}{name={\textsf{\textcolor{red}{\term{PCLK} (Pixel Clock)}}},sort={00855},description={Output of the \term{OV7670} clocking its data bus, one cycle per transferred byte. \cite{OV7670-DS}}}
\newglossaryentry{picocom}{name={\textsf{\textcolor{red}{\texttt{picocom}}}},sort={0086},description={Minimal terminal emulator for serial lines, used to reach the console of the board. \citegloss{PICOCOM-REPO}}}
\newglossaryentry{pio}{name={\textsf{\textcolor{red}{\term{PIO} (Parallel Input/Output)}}},sort={0087},description={\term{Intel}/\term{Altera} \term{IP} core that exposes a set of \term{FPGA} pins as a memory-mapped register, letting software drive or read them directly; used here, for example, to light the on-board \term{LED}s. \cite{EP-IPUG}}}
\newglossaryentry{platformdesigner}{name={\textsf{\textcolor{red}{\term{Platform Designer} (\term{Qsys})}}},sort={0088},description={\term{Quartus Prime} tool for assembling an embedded system from \term{IP} cores and generating the interconnect between them. \term{Qsys} is its former name, still visible in the \texttt{.qsys} file extension. \cite{ALTERA-QUARTUS}}}
\newglossaryentry{pollingmode}{name={\textsf{polling mode}},sort={0089},description={Operating mode in which software repeatedly reads a status register to detect an event, the opposite strategy to interrupt-driven notification.}}
\newglossaryentry{prefetcher}{name={\textsf{\textcolor{red}{prefetcher}}},sort={0090},description={Hardware block inside the \term{mSGDMA} core that reads a chain of committed descriptors from a dedicated descriptor memory area and hands each one to the dispatcher on its own, so software only has to point it at the first descriptor rather than handing them over one at a time. In this work that memory is on-chip; \prjname{FPGAsteel} enables the prefetcher on every \term{mSGDMA} instance, while \prjname{FPGAlix} leaves it disabled and manages descriptors from software instead. \cite{EP-IPUG}}}
\newglossaryentry{prompt}{name={\textsf{prompt}},sort={0091},description={String printed by a program to signal that it is ready and waiting for input. Shells and bootloaders are common examples, but any interactive program can print one, such as this project's \texttt{vision\_core} or its debug scripts.}}
\newglossaryentry{putty}{name={\textsf{\textcolor{red}{\term{PuTTY}}}},sort={0092},description={\term{Windows} terminal emulator supporting serial and \term{SSH} connections. \citegloss{PUTTY-SITE}}}
\newglossaryentry{quartusprime}{name={\textsf{\textcolor{red}{\term{Quartus Prime}}}},sort={0093},description={\term{Intel}/\term{Altera} design suite for \term{FPGA} synthesis, placement and routing, timing analysis and device programming. \cite{ALTERA-QUARTUS}}}
\newglossaryentry{registerinterface}{name={\textsf{register interface}},sort={0094},description={Set of memory-mapped registers that a hardware core exposes to software for configuration, command and status reporting.}}
\newglossaryentry{resetmanager}{name={\textsf{\textcolor{red}{Reset Manager}}},sort={0095},description={\term{HPS} peripheral through which software asserts and releases the reset of individual \term{HPS} components, including each \term{Cortex-A9} core; used in \prjname{FPGAsteel} to bring \term{CPU1} out of reset. \cite{CV-TRM}}}
\newglossaryentry{rgb444}{name={\textsf{\textcolor{red}{\term{RGB444}}}},sort={0096},description={Twelve-bit pixel format with four bits each for red, green and blue, packed into a sixteen-bit word with four unused bits. \cite{OV7670-DS}}}
\newglossaryentry{rgb555}{name={\textsf{\textcolor{red}{\term{RGB555}}}},sort={0097},description={Fifteen-bit pixel format with five bits each for red, green and blue, packed into a sixteen-bit word with one unused bit. \cite{OV7670-DS}}}
\newglossaryentry{rgb565}{name={\textsf{\textcolor{red}{\term{RGB565}}}},sort={0098},description={Sixteen-bit pixel format with five bits for red, six for green and five for blue, exactly filling a sixteen-bit word; it is the format produced by the acquisition path of this project. \cite{OV7670-DS}}}
\newglossaryentry{rootfilesystem}{name={\textsf{root filesystem}},sort={0099},description={Filesystem mounted at \texttt{/} that contains the userland of the system: the init process, the shell, the libraries, the configuration files and the applications. The kernel mounts it at the end of boot; without it the kernel has nothing to execute.}}
\newglossaryentry{rtsp}{name={\textsf{\textcolor{red}{\term{RTSP} (Real Time Streaming Protocol)}}},sort={0100},description={Control protocol used to establish and steer media streams between a server and its clients. \citegloss{RFC2326-RTSP}}}
\newglossaryentry{scp}{name={\textsf{\textcolor{red}{\texttt{scp}}}},sort={0101},description={\term{OpenSSH} secure file copy: command-line tool that copies files between host and target over an encrypted \term{SSH} connection, using a syntax similar to \texttt{cp} but with a \texttt{user@host:path} form to name the remote side. Used throughout this work to deploy compiled binaries and debug scripts onto the board. \citegloss{SCP-MANPAGE}}}
\newglossaryentry{scu}{name={\textsf{\textcolor{red}{\term{SCU} (Snoop Control Unit)}}},sort={0102},description={\term{Cortex-A9} block that keeps the per-core \term{L1} caches coherent with each other, snooping one core's cache on behalf of the other; a core joins its coherency domain by setting the \term{ACTLR.SMP} bit. \cite{CV-TRM}}}
\newglossaryentry{serialconsole}{name={\textsf{serial console}},sort={0103},description={Text console exposed over a \term{UART}. On the \term{DE1-SoC} it is reached through the on-board \term{USB}-to-\term{UART} bridge, so it appears on the host as a serial port over \term{USB}; it carries, for example, the bootloader and kernel boot log, or, in \prjname{FPGAsteel}, whatever a bare-metal application prints directly.}}
\newglossaryentry{shell}{name={\textsf{shell}},sort={0104},description={Program that reads the commands typed by the user, interprets them and launches the corresponding executables, returning their output. It is the usual way of driving a \term{UNIX}-like system without a graphical interface, and it is also a programming language in its own right, since the same commands can be collected into a script and run unattended. \term{Bash} is the shell of most \term{Linux} distributions; on the board \term{BusyBox} provides a lighter one. Beyond the operating system, the term also denotes the interactive command interpreter offered by other software, such as the one \term{U-Boot} presents when the automatic boot is interrupted.}}
\newglossaryentry{signaltap}{name={\textsf{\textcolor{red}{\term{SignalTap}}}},sort={0105},description={Embedded logic analyzer of \term{Quartus Prime}, which captures internal \term{FPGA} signals at run time and uploads the acquired traces to the host over \term{JTAG}. \citegloss{QUARTUS-DEBUG-TOOLS}}}
\newglossaryentry{sigterm}{name={\textsf{\textcolor{red}{\term{SIGTERM} and \term{SIGINT}}}},sort={0106},description={\term{POSIX} signals that request a process to terminate: \term{SIGTERM} is the standard termination request, \term{SIGINT} the one raised from the terminal with \term{Ctrl-C}. \citegloss{GLIBC-TERMINATION-SIGNALS}}}
\newglossaryentry{singleton}{name={\textsf{\textcolor{red}{singleton}}},sort={0107},description={Design pattern, implemented here in \term{C++}, that restricts a class to a single instance and provides a global point of access to it. \citegloss{GREGOIRE-PROFESSIONAL-CPP}}}
\newglossaryentry{smp}{name={\textsf{\term{SMP} (Symmetric Multiprocessing)}},sort={0108},description={Architecture in which several identical \term{CPU} cores share the same memory, each capable of running any task.}}
\newglossaryentry{sobeledgedetection}{name={\textsf{\textcolor{red}{Sobel edge detection}}},sort={0109},description={Edge detection operator that approximates the gradient of the image intensity along the horizontal and vertical directions; edges are then extracted wherever the resulting gradient magnitude is large, i.e.\ where intensity changes sharply. \citegloss{GONZALEZ-WOODS-DIP}}}
\newglossaryentry{socfpga}{name={\textsf{\textcolor{red}{\term{SoC FPGA}}}},sort={0110},description={Device that integrates a \term{CPU} and an \term{FPGA} fabric on a single die: software runs on the \term{CPU}, custom digital hardware on the \term{FPGA} fabric, combining the flexibility of the former with the acceleration of the latter. \citegloss{MICROCHIP-SOCFPGA}}}
\newglossaryentry{socsystemqsys}{name={\textsf{\texttt{soc\_system.qsys}}},sort={0111},description={\term{Platform Designer} project file describing the complete hardware system assembled for this work, including every \term{IP} core it instantiates and the interconnect between them.}}
\newglossaryentry{sop}{name={\textsf{\textcolor{red}{\term{SOP} and \term{EOP} (\texttt{st\_sop}, \texttt{st\_eop})}}},sort={0112},description={Start-of-packet and end-of-packet markers of \term{Avalon-ST}, which delimit a packet on the stream. In this design a packet is one video frame. \cite{AVALON-SPEC}}}
\newglossaryentry{spl}{name={\textsf{\textcolor{red}{\term{SPL} (Secondary Program Loader)}}},sort={0113},description={First stage of \term{U-Boot}, loaded by the boot \term{ROM} into the on-chip \term{RAM} of the \term{HPS}. It initializes the \term{SDRAM} controller and the pin multiplexing, then loads the full \term{U-Boot} image. \cite{SUB-UBOOT,CV-TRM}}}
\newglossaryentry{ssh}{name={\textsf{\textcolor{red}{\term{SSH} (Secure Shell)}}},sort={0114},description={Protocol for encrypted remote login and file transfer over an untrusted network. \citegloss{RFC4251-SSH}}}
\newglossaryentry{su}{name={\textsf{\textcolor{red}{\texttt{su}}}},sort={0115},description={\term{Linux} command that starts a shell as another user; invoked without arguments it switches to root. \citegloss{SU-MANPAGE}}}
\newglossaryentry{symbol}{name={\textsf{\textcolor{red}{symbol}}},sort={0116},description={Elementary data unit of an \term{Avalon-ST} link, whose width is a property of the interface and is chosen so that it corresponds to a meaningful quantity: eight bits, one byte of pixel data, in the camera stream. A beat carries an integer number of symbols, so the symbol width, together with the number of symbols per beat, fixes how much data the link moves per beat. \cite{AVALON-SPEC}}}
\newglossaryentry{tlb}{name={\textsf{\textcolor{red}{\term{TLB} (Translation Lookaside Buffer)}}},sort={0117},description={Per-core cache of recently used \term{MMU} page table entries, avoiding a memory lookup on every address translation; each \term{Cortex-A9} core has its own, even when both point at the same page table. \citegloss{BRYANT-OHALLARON-CSAPP}}}
\newglossaryentry{uboot}{name={\textsf{\textcolor{red}{\term{U-Boot}}}},sort={0118},description={Bootloader widely used in embedded systems. It configures the hardware, offers an interactive command shell and loads the kernel or the bare-metal application. \cite{SUB-UBOOT}}}
\newglossaryentry{uart0}{name={\textsf{\textcolor{red}{\term{UART0}}}},sort={0119},description={First serial port of the \term{HPS}, routed on the \term{DE1-SoC} to the \term{USB}-to-\term{UART} bridge and exposed on a \term{micro-USB} connector; it carries the serial console. \cite{CV-TRM,DE1-SCH}}}
\newglossaryentry{udev}{name={\textsf{\textcolor{red}{\texttt{udev}}}},sort={0120},description={Device manager of \term{Linux}, which creates the nodes under \texttt{/dev} and applies the rules that govern them, for instance the ones granting user access to the \term{USB-Blaster}. \citegloss{UDEV-MANPAGE}}}
\newglossaryentry{usbblaster}{name={\textsf{\textcolor{red}{\term{USB-Blaster} and \term{USB-Blaster II}}}},sort={0121},description={\term{Intel}/\term{Altera} \term{USB}-to-\term{JTAG} programming and debugging interface. The \term{DE1-SoC} integrates a \term{USB-Blaster II} directly on the board. \citegloss{ALTERA-USB-BLASTER-UG}}}
\newglossaryentry{v4l2}{name={\textsf{\textcolor{red}{\term{V4L2} (Video4Linux2)}}},sort={0122},description={\term{Linux} kernel \term{API} for video capture devices, defining a standard interface for format negotiation, buffer queues and streaming; used here by \texttt{\seqsplit{camera\_driver\_2}}. \citegloss{KERNEL-V4L2-DOC}}}
\newglossaryentry{vfp}{name={\textsf{\textcolor{red}{\term{VFP} (Vector Floating-Point)}}},sort={0123},description={\term{ARM} coprocessor that executes floating-point instructions; on the \term{Cortex-A9} its access must be explicitly enabled in software before use, together with \term{NEON}. \citegloss{ARM-DDI-0409I}}}
\newglossaryentry{vga}{name={\textsf{\term{VGA}}},sort={0124},description={Used here in two senses: the 640$\times$480 resolution at which the sensor is operated, and the analog display interface driven by the output path of \prjname{FPGAsteel}.}}
\newglossaryentry{vhdl}{name={\textsf{\textcolor{red}{\term{VHDL}}}},sort={0125},description={Hardware description language used to specify the behavior and the structure of the digital logic synthesized into the \term{FPGA}. \cite{PEDRONI-VHDL}}}
\newglossaryentry{visualstudiocode}{name={\textsf{\textcolor{red}{\term{Visual Studio Code}}}},sort={0126},description={Source code editor developed by \term{Microsoft} and distributed free of charge, used here to write both the \term{VHDL} sources and the \term{C/C++} application code. \citegloss{VSCODE-SITE}}}
\newglossaryentry{vm}{name={\textsf{\textcolor{red}{\term{VM} (Virtual Machine)}}},sort={0127},description={Software emulation of a complete computer, running its own operating system on top of a host, through a software called hypervisor. \citegloss{VMWARE-VM}}}
\newglossaryentry{vsync}{name={\textsf{\textcolor{red}{\term{VSYNC}}}},sort={0128},description={Output of the \term{OV7670} asserted once per frame, which marks the frame boundary in the acquisition path. \cite{OV7670-DS}}}
\newglossaryentry{vtd}{name={\textsf{\term{VT-d}, \term{AMD-Vi} and \term{IOMMU}}},sort={0129},description={\term{Intel}'s \term{VT-d} and \term{AMD}'s \term{AMD-Vi} are the respective \term{x86\_64} hardware virtualization extensions for \term{I/O}. Both provide an \term{IOMMU} (\term{I/O} Memory Management Unit), which translates the addresses issued by devices and enables reliable \term{PCI(e)} device pass-through to a virtual machine.}}
\newglossaryentry{xpackbuilds}{name={\textsf{\textcolor{red}{xPack builds}}},sort={0130},description={Distribution of prebuilt, version-pinned development tools, such as toolchains and \term{\seqsplit{OpenOCD}}. \cite{XPACK-ARM-GCC,XPACK-OPENOCD}}}
\newglossaryentry{yuv422}{name={\textsf{\textcolor{red}{\term{YUV} 4:2:2}}},sort={0131},description={Pixel format that separates the luma component from two chroma components and samples the chroma at half the horizontal resolution; it is one of the output formats supported by the \term{OV7670}. \cite{OV7670-DS}}}
